\minitopic{研读实验室：从直觉到可复核计算}\index{条件概率}\index{贝叶斯因子}\index{校准}\index{期望效用} 形式认识论的价值不在于用公式替代判断，而在于公开判断结构。一个后验概率必须说明先验、似然与条件独立；一个决策必须说明可能结果和效用；一个专家汇总必须说明权重。读者可以不同意输入，却能复算结论并定位分歧。 概率记号首先防止方向混淆。$P(D\mid +)$表示检测阳性后患病概率，$P(+\mid D)$表示患病者检测阳性的概率。后者是灵敏度，不能直接当作前者。贝叶斯定理把两者连接：
\[
P(D\mid +)=\frac{P(+\mid D)P(D)}
{P(+\mid D)P(D)+P(+\mid \neg D)P(\neg D)}.
\]
分母保留了真阳性和假阳性两种到达“阳性”的路径。 设患病率为1\%，灵敏度90\%，假阳性率5\%。在一万人中约有100名患者，其中90人阳性；9900名非患者中约495人阳性。所有585名阳性者中只有90名患病，因此
\[
P(D\mid +)=\frac{90}{90+495}\approx15.4\%.
\]
“检测准确率很高”若不区分错误类型和基础率，容易产生远高于15.4\%的直觉。

\minitopic{赔率形式与贝叶斯因子}概率$p$可转成赔率$p/(1-p)$。上述先验赔率是$0.01/0.99\approx1:99$。阳性的贝叶斯因子为
\[
BF_+=\frac{P(+\mid D)}{P(+\mid\neg D)}=\frac{0.90}{0.05}=18.
\]
后验赔率等于先验赔率乘以18，约为$18:99$，再转回概率即15.4\%。赔率形式清楚显示证据贡献与先验分离。 若第二次检测在给定患病状态后与第一次条件独立，两个阳性的总贝叶斯因子为$18^2=324$，后验赔率约$324:99$，概率约76.6\%。若两次都使用同一污染样本，条件独立不成立；简单平方会把共享错误当作两份支持。 贝叶斯因子衡量证据区分两个假设的能力，不等于某假设的最终概率。同一阳性对低基础率和高基础率人群给出相同似然比，却产生不同后验。报告“证据强18倍”时，应说明比较的是哪两个假设以及参数如何估计。

\minitopic{校准、分辨与评分}一个预测者若在所有报70\%的事件中约七成发生，就在该档位校准良好。校准不要求每个70\%预测都成真；单次失败不能证明概率错。长期检验应把相近概率分组，比较预测频率与实际频率，并报告样本量和不确定区间。 校准也不等于信息丰富。总说“明天下雨概率等于当地长期频率”的预测可以校准，却不能区分具体天气。分辨能力要求在会下雨的日子给更高概率、不会下雨的日子给更低概率。好系统同时追求校准和分辨。 恰当评分规则奖励诚实概率报告。以Brier分数为例，二元事件结果$o$取0或1，预测为$p$，损失为$(p-o)^2$。报0.9而事件未发生的损失很大，报0.6的损失较小；长期上，按真实置信报告可最小化期望损失。评分规则评价置信质量，不自动判定主体是否拥有知识\parencite{joyce1998}。 预测系统还需防止选择性展示。只公布成功的高置信预测、删除失败记录，会制造虚假校准。应在结果发生前锁定预测、保存全部版本并预先规定评估窗口。形式透明依赖记录制度。

\minitopic{先验、参照类与模型选择}先验不是纯主观偏好。它可来自基础率、先前研究、机制知识或对称性原则。但不同参照类给出不同基础率：某病在全国、某年龄组、某基因携带者中的比例不同。选择应以与目标个体因果和抽样相关的特征为依据，不能事后挑选得到所需结论的群体。 先验也不能包含当前证据两次。若基础率研究已使用这名患者的检测结果，再把同一结果作为新似然更新，就会重复计算。历史数据与当前证据的边界在大数据系统中尤其难，需要记录训练集、验证集和个案输入。 复杂模型通常能更好拟合已有数据，却可能过拟合噪声。模型比较要考察新数据预测、参数惩罚与稳健性，而非只看训练拟合。旧证据问题说明，一个理论若在看到证据后才调参解释它，解释成功的确认价值低于事前给出新预测。 先验分歧不必导致永久相对主义。若不同主体使用非零且不过度极端的先验，持续获得具有区分力的共同证据时，后验往往趋近。但若双方模型赋予证据完全不同似然、或只接触各自选择的信息，分歧可持续。形式化帮助说明分歧在哪个输入，而非保证自动共识。

\minitopic{从概率到行动}概率不包含损失。即使患病概率只有5\%，若漏诊致命而进一步检查成本很低，行动上仍应检查；即使某市场上涨概率60\%，潜在损失极大也未必投资。期望效用把每项行动在各状态的效用按概率加权\parencite{jeffrey1983}。 设治疗在患病时带来$+100$，未患病却治疗为$-20$；不治疗在患病时为$-200$，未患病为0。若患病概率为$q$，治疗期望效用为$100q-20(1-q)=120q-20$，不治疗为$-200q$。治疗优于不治疗当$320q>20$，即$q>6.25\%$。行动阈值来自效用，不是知识定义。 效用数字本身需要伦理和经验论证。不同人的副作用负担、风险态度和权利限制不能随意压成一个专家数字。法律定罪尤其不能只按社会总效用选择概率阈值，因为程序正义和个体权利构成约束。决策模型公开权衡，却不能替我们决定所有价值。 信息也有价值。若一次额外检测可能显著改变行动，而且成本低，其期望信息价值高；若无论结果如何都采取同一行动，继续检测可能无益。信息价值解释何时应暂停行动搜集证据，也解释为什么高风险不总意味着无限调查。

\minitopic{比较视角：墨家“三表”与形式证据}后期文献传统常以墨家“三表”讨论言说应考察古代依据、百姓耳目经验和实施效果；具体篇章、版本与解释存在争议，不能把它直接写成贝叶斯定理\parencite{graham1989}。其可比较之处是，主张要同时受传统记录、公共经验和实践后果约束，而非只凭权威或辞令。 三类检验之间没有现成数值权重，也未给出条件概率公理。当代形式工具可以追问：材料是否独立、历史记录的基础率怎样、实施效果是否由目标主张造成；墨家问题意识则提醒，证据评价与公共可检验性、治理后果密切相关。 把古代“数”或“法”见到就译成统计模型，是时代错置；反过来，认为没有现代符号便没有严谨推理，也忽略论证规则可以用非数学方式表达。比较应在证据来源、公开检验与行动反馈三个问题层面进行，不宣称形式同一。

\minitopic{综合案例：大学录取模型}学校模型给申请者70\%的“学业成功概率”。首先要问事件定义：一年不退学、四年毕业还是绩点达标？训练人群和当前申请者是否同分布？概率对不同群体是否分别校准？总校准可能掩盖交叉群体的系统误差。 模型使用高中成绩、课程资源和历史录取数据。若历史录取本身受机会不平等影响，预测可能准确再现现状，却不能直接回答谁应被录取。预测命题与规范决策不同：成功概率是证据输入，教育潜力、补偿目标和权利约束还需单独论证。 招生员又读取推荐信，而推荐信内容已被模型文本特征使用。如果人工判断再次把同一内容当独立证据，就发生重复计算。应标记哪些信息已经进入模型，哪些是新证据，并用盲审测试人工增量是否真正提高校准。 最终报告不应只给单分数。它应给出事件定义、置信区间、适用人群、主要特征、已知失效模式和对阈值附近个案的复核流程。形式认识论的目标不是让数字获得神秘权威，而是让推断与决定可追踪、可反驳。

\begin{tcolorbox}[breakable,colback=white,colframe=blue!30!black,title={本章研读任务},fonttitle=\bfseries]
任选一个检测或预测，使用一万人自然频率计算后验，再以赔率和贝叶斯因子复算；检查两份证据是否条件独立；设计校准表和评分方式；最后给出一个效用矩阵，明确区分概率结论、知识判断与行动阈值。
\end{tcolorbox}
